\documentclass[11pt]{article}
% Change "review" to "final" to generate the final (sometimes called camera-ready) version.
% Change to "preprint" to generate a non-anonymous version with page numbers.
\usepackage[]{acl}
\usepackage{times}
\usepackage{latexsym}

% For proper rendering and hyphenation of words containing Latin characters (including in bib files)
\usepackage[T1]{fontenc}
% For Vietnamese characters
% \usepackage[T5]{fontenc}
% See https://www.latex-project.org/help/documentation/encguide.pdf for other character sets

% This assumes your files are encoded as UTF8
\usepackage[utf8]{inputenc}
\usepackage{microtype}
\usepackage{inconsolata}
\usepackage{graphicx}

% If the title and author information does not fit in the area allocated, uncomment the following
%
%\setlength\titlebox{<dim>}
%
% and set <dim> to something 5cm or larger.

\title{Group 10 Final Report:\\Idiom Classification}


\author{Rebecca Di Filippo, Leo Vugert, Timothy Pokani, Claire Nielsen \\
  \texttt{\{difilir,vugertl,pokanait,nielsc2\}@mcmaster.ca} }

%\author{
%  \textbf{First Author\textsuperscript{1}},
%  \textbf{Second Author\textsuperscript{1,2}},
%  \textbf{Third T. Author\textsuperscript{1}},
%  \textbf{Fourth Author\textsuperscript{1}},
%\\
%  \textbf{Fifth Author\textsuperscript{1,2}},
%  \textbf{Sixth Author\textsuperscript{1}},
%  \textbf{Seventh Author\textsuperscript{1}},
%  \textbf{Eighth Author \textsuperscript{1,2,3,4}},
%\\
%  \textbf{Ninth Author\textsuperscript{1}},
%  \textbf{Tenth Author\textsuperscript{1}},
%  \textbf{Eleventh E. Author\textsuperscript{1,2,3,4,5}},
%  \textbf{Twelfth Author\textsuperscript{1}},
%\\
%  \textbf{Thirteenth Author\textsuperscript{3}},
%  \textbf{Fourteenth F. Author\textsuperscript{2,4}},
%  \textbf{Fifteenth Author\textsuperscript{1}},
%  \textbf{Sixteenth Author\textsuperscript{1}},
%\\
%  \textbf{Seventeenth S. Author\textsuperscript{4,5}},
%  \textbf{Eighteenth Author\textsuperscript{3,4}},
%  \textbf{Nineteenth N. Author\textsuperscript{2,5}},
%  \textbf{Twentieth Author\textsuperscript{1}}
%\\
%\\
%  \textsuperscript{1}Affiliation 1,
%  \textsuperscript{2}Affiliation 2,
%  \textsuperscript{3}Affiliation 3,
%  \textsuperscript{4}Affiliation 4,
%  \textsuperscript{5}Affiliation 5
%\\
%  \small{
%    \textbf{Correspondence:} \href{mailto:email@domain}{email@domain}
%  }
%}

\begin{document}
\maketitle
% \begin{abstract}
% \end{abstract}

\section{Introduction}

For this project, the team has created software that identifies and classifies idiomatic expressions and literary devices in English text. Specifically, this task
focuses on detecting and labelling the type of figurative language present in short text segments, including: metaphors, similes, personification, as well as a few others. \\

Understanding idiomatic and figurative language is a fundamental challenge in natural language processing because the intended meaning often cannot be
derived directly from the literal meanings of the individual words. While humans can usually recognize these expressions using context and world knowledge,
computational models struggle due to ambiguity, cultural conventions, and the subtle differences between similar devices (for example, discerning a metaphor vs. an idiom).\\

This task is significant because idiomatic language appears frequently in real world applications of natural language processing such as social media analysis, educational tools, and conversational agents. Misinterpreting these expressions can lead to incorrect or misleading system responses. Unlike standard sentiment analysis, which focuses primarily on emotional polarity, this task requires the software to perform semantic disambiguation and contextual reasoning. This goes beyond surface-level text classification, and requires the model to navigate linguistic ambiguity where the intended meaning cannot be derived directly from the literal meanings of individual words.\\

\section{Related Work}

The team found their dataset attached to a study done in 2022, which performed the same task. This study was incredibly helpful as a benchmark, as it tested for the same categories as this task. This was the first indication that BERT would perform quite well, as in this study it outperformed both Multinomial Naive-Bayes and Support Vector Machine models \citet{adewumi2022potentialidiomaticexpressionpieenglish}. The BERT resultant F1-Score was 0.948, which this team would later surpass. \\

A study performed in 2009 examined the results of a combined classifier that uses both supervised and unsupervised models to examine literal vs. non-literal use of idiomatic expressions. This study found that the combined classifier was more accurate than the unsupervised classifier alone and the majority baseline; the combined classifier resulted in an error reduction of over 35\% \cite{li-sporleder-2009-classifier}. \\

Research by \citet{11251805} highlighted the success of the BERT model when applied to a text to analyse literary devices. Models such as GloVe and Word2Vec treat words as separate entities, while BERT takes in the surrounding context, leading to much better performance analysing literary devices. This paper reported an F1-Score of 0.85 resulting from applying BERT to detect metaphors on a corpus build from Project Gutenberg and other public literary anthologies. \\

A similar paper published in 2024 tested the BERT model on text classification, sentiment analysis, and and question answering. This paper found that BERT yielded an F1-Score on text classification of 0.915, a score on sentiment analysis of 0.885, and a score on question answering of 0.865. These results were all higher than their counterparts performed by recurrent neural networks (RNNs), supporting the conclusion that the BERT model would outperform RNN methods, at the cost of much longer training time and computational load \cite{10918499}. \\

A study was done by \citet{wu2024researchapplicationdeeplearningbased} that compared BERT and DistilBERT while performing sentiment analysis. It explains that by updating parameter weights used on the BERT model, DistilBERT can perform better on sentence classification when using a full set of parameters. \\

\section{Dataset}

\subsection{Source and Collection}
The dataset utilized for this project is the Potential Idiomatic Expression (PIE)-English Corpus, which is a labeled dataset designed for the classification of idiomatic expressions and literary devices. The data is publicly available via a GitHub repository and is derived primarily from the British National Corpus (BNC), which accounts for 96.9\% of the samples, and UK Web Pages (UKWaC), which accounts for the remaining 3.1\%. The original corpus was extracted by near-native English speakers and labeled by independent annotators to ensure accuracy. 

\subsection{Dataset Properties and Labels}
The task is defined as a single-label classification problem where each text segment is mapped to one of ten primary literary device or idiomatic categories. The ten classes include: Euphemism, Hyperbole, Irony, Literal, Metaphor, Oxymoron, Paradox, Parallelism, Personification, and Simile. 

A critical characteristic of this dataset is the significant class imbalance. As shown in Table \ref{tab:label_dist}, the distribution is heavily skewed toward Metaphors, while some categories, such as Irony and Oxymoron, represent less than 1\% of the total data points. This distribution presents a notable challenge for model training, as the "long-tail" labels have very low frequency.

\begin{table}[h]
  \centering
  \caption{Label Distribution in the PIE-English Corpus}
  \label{tab:label_dist}
  \resizebox{\columnwidth}{!}{% 
  \begin{tabular}{|l|r|l|r|}
  \hline
  \textbf{Label} & \textbf{Frequency} & \textbf{Label} & \textbf{Frequency} \\ \hline
  Metaphor & 72.7\% & Oxymoron & 0.24\% \\ \hline
  Euphemism & 11.82\% & Paradox & 0.56\% \\ \hline
  Simile & 6.11\% & Hyperbole & 0.24\% \\ \hline
  Literal & 5.65\% & Parallelism & 0.32\% \\ \hline
  Personification & 2.22\% & Irony & 0.16\% \\ \hline
  \end{tabular}%
  }
  \end{table}
\subsection{Preprocessing}
The original source data was provided in a token-per-row CSV format. To make the data suitable for sentence-level classification and manual annotation, we implemented a Python script to perform the following preprocessing steps:
\begin{itemize}
    \item \textbf{Sentence Reconstruction:} Tokens were joined to reconstruct full sentences from the tokenized rows.
    \item \textbf{Normalization:} All reconstructed sentences were converted to lowercase.
    \item \textbf{Random Sampling:} From the total 20,174 samples, we randomly sampled 1,200 data points to create a balanced mixture for our team's manual annotation phase.
\end{itemize}

\subsection{Annotation Procedure and Reliability}
Our team followed a structured annotation protocol using a shared spreadsheet interface. Each of the four team members independently annotated 300 unique data points, resulting in 1,200 total primary annotations. To measure Inter-Annotator Agreement (IAA), each member also re-annotated a subset of 45 data points originally labeled by a teammate.

The consistency of our annotations was evaluated using Cohen's Kappa for pairwise comparisons and Krippendorff's Alpha for overall group reliability. Our pairwise Cohen's Kappa scores ranged from 0.6293 to 0.8990, representing substantial to great agreement. The final Krippendorff's Alpha was 0.8790, which actually exceeded some individual pairwise scores, suggesting our group was highly aligned with the original corpus's "Gold Standard" labels

\subsection{Data Observations}
During the annotation process, we identified several interesting linguistic nuances. Some sentences presented "edge cases" where literal and metaphorical meanings overlapped, such as \textit{"the worry at the back of her mind,"} which is metaphorical, appearing alongside the literal phrase \textit{"she sat down beside her patient"}. Furthermore, because the corpus is British-centric, we encountered regional idioms such as \textit{"not by a long chalk,"} which is the UK equivalent of the North American \textit{"not by a long shot"}. These cultural variations and internal sentence ambiguities highlight the difficulty of the classification task for automated models.


\section{Features}

\subsection{Model Inputs and Representation}
The primary input for our models consists of reconstructed English sentences derived from the token-per-row CSV format of the PIE-English Corpus. During the preprocessing phase, these tokens are joined to form complete grammatical units, ensuring the models have access to the full sentential context necessary for disambiguating figurative language. To reduce the complexity of the input space and handle vocabulary sparsity, all text is normalized to lowercase before being fed into the embedding layers.

\subsection{Representation Learning and Embeddings}
Our team utilizes a combination of learned and
 pretrained embeddings. For our transformer-based models, 
 we perform representation learning where the model learns 
 task-specific embeddings through fine-tuning.This approach
  is vital because the meaning of an idiom often cannot 
  be derived from individual word embeddings alone. 
  By fine-tuning, the model develops a representation 
  that prioritizes the semantic relationship between words
   in a figurative context.

\subsection{Feature Engineering}
In addition to automated representation learning, we incorporate manual feature engineering to assist the models in identifying structural cues for low-frequency classes. Key features include:
\begin{itemize}
    \item \textbf{Part-of-Speech (POS) Tagging:} We expect POS tags to be highly useful for identifying idioms that follow rigid grammatical patterns, such as the noun-preposition-noun structure.
    \item \textbf{Syntactic Pattern Matching:} For categories like similes, we explicitly look for comparative markers such as "as" or "like." For oxymorons, we look for the juxtaposition of contradictory adjectives and nouns.
    \item \textbf{Contextual Tone:} We analyze the overall tone and sentiment of the sentence to help distinguish between literal statements and irony, as irony often involves a sarcastic or contradictory emotional polarity compared to the literal text.
\end{itemize}

\subsection{Feature Selection and Data Augmentation}
Due to the extreme class imbalance, where metaphors represent over 72\% of the dataset - we decided to focus on feature selection that highlights devices like paradoxes and ironies, which make up less of the data. Rather than removing data, we utilize the entire dataset to ensure the models are exposed to as many minority-class examples as possible. We also consider feature augmentation by providing the model with binary flags for the presence of specific keywords known to trigger figurative labels, assisting the model in learning patterns associated with low-frequency classes.
\section{Implementation}

Describe your model and implementation here. Refer to item 4. This may take around a page.
TIM/LEO
\section{Results and Evaluation}

What are your results? What did you implement and what did your classmates implement? What are your baselines? Refer to items 5 and 6. This may take around 0.5-1 pages.
BECCA
\section{Comparison to Others}

What are your results? What did you implement and what did your classmates implement? What are your baselines? Refer to items 5 and 6. This may take around 0.5-1 pages.
CLAIRE
\section{Error Analysis}

What kind of errors are your models making? This doesn't have to be super complicated. Confusion matrices are great. Qualitative observations are also great. This means looking at many cases where your model is wrong (wrong class or high error) and looking for patterns yourself. What about the text might be giving it difficulty? Qualitative analysis is underappreciated and can often give you great insight into issues with your model. Examples are useful here. Refer to item 7. This may be around a half page.
TIM
\section*{Conclusions}

Summarize your work briefly. Reflect on your implementation and areas for improvement. What do you suggest people work on next (mentioned also in item 7)? This should likely be less than 0.25 pages.
BECCA

\section{Team Contributions} The members of Team 10 each took on specialized roles to ensure the successful execution of both the technical code and documentation aspects of the project.
\begin{itemize}
  \item \textbf{Rebecca Di Filippo}: Created repository and main task was helping create the final report. Also assisted choosing coding methods for models.
  \item \textbf{Claire Nielsen}:  Led the documentation efforts, including writing the final report and ensuring that all sections were cohesive and well-structured. 
  \item \textbf{Tim Pokanai}: Implemented 2 of the 4 models, and helped right the implementation and error analysis section.
  \item \textbf{Leo Vugert}: Implemented 2 of the 4 models, and helped right the implementation and error analysis section.
  \end{itemize}
%\bibliography{final_report/custom}

\end{document}
